% AUTORIDAD DE FUENTE: Exposé IX, §14.3, líneas 4627--4723;
% índices de escaneo 525--529, folios 514--518. Las capas francesas D/C
% no presentan diferencias materiales. Seis errores tipológicos o de
% referencia inequívocos se conservan reversiblemente mediante \sisi y
% se han comunicado al custodio raíz. La coma espuria del testigo inglés
% en \underline m^{,n+1} se rechazó contra ambas capas francesas, que
% coinciden en la lectura correcta \underline m^{n+1}.

\par\medskip
\noindent
14.3. \underline{Aplicación a una caracterización aritmética de}
$i:\underline M_K\longrightarrow G_K$.

La caracterización que tenemos en mente se basa en la observación de
que, para todo número primo $\ell$, la extensión (14.2.2) de
$\underline M_K\otimes\mathbb Z_\ell$ por $T_\ell(G_K)$ deducida de
$i$ se conoce por medios puramente algebraicos. En particular, se
conoce la imagen de $i$ por el homomorfismo canónico

\begin{equation}\tag{14.3.1}
 \operatorname{Hom}(\underline M_K,G_K)\longrightarrow
 \prod_\ell\operatorname{Ext}^1
 \bigl(\underline M_K\otimes\mathbb Z_\ell,T_\ell(G_K)\bigr)
 \qquad .
\end{equation}

\noindent
En el caso (al que siempre puede reducirse uno por descenso) en que
$\underline M_K$ es constante, de valor $M$, el homomorfismo precedente
puede interpretarse como

\begin{equation}\tag{14.3.2}
 \operatorname{Hom}\bigl(M,G_K(K)\bigr)\longrightarrow
 \prod_\ell
 \sisi{\operatorname{Hom}\bigl(M,H^1(K,T_\ell(G_K))}%
 {\operatorname{Hom}\bigl(M,H^1(K,T_\ell(G_K))\bigr)},
\end{equation}

\noindent
deducido del homomorfismo canónico

\begin{equation}\tag{14.3.3}
 G_K(K)\longrightarrow\prod_\ell H^1(K,T_\ell(G_K))
 \qquad ,
\end{equation}

\noindent
que proviene de los homomorfismos de conexión de las sucesiones exactas
de Kummer

\begin{equation*}
 0\longrightarrow{}_nG_K\longrightarrow G_K
 \xrightarrow{\ n\ }G_K\longrightarrow0\qquad .
\end{equation*}

%%%%%%%%%%  scan idx 526  |  folio 515  %%%%%%%%%%

(El homomorfismo (14.3.1) podría, de hecho, interpretarse siempre como
el homomorfismo de conexión análogo a (14.3.3) para el grupo torcido
$\check{\underline M}_K\otimes G_K=
\underline{\operatorname{Hom}}(\underline M_K,G_K)$.) El núcleo de
(14.3.3) está formado manifiestamente por los elementos de $G_K(K)$ que
son \underline{divisibles por} $n$ \underline{para todo entero}
$n\geq1$, o, como diremos, que son \underline{infinitamente
divisibles}. Esto determina ya $i$ módulo un elemento de
$\operatorname{Hom}(M,g)$, donde $g$ es el subgrupo de $G_K(K)$ formado
por sus elementos infinitamente divisibles:

\begin{equation*}
 g=\bigcap_{n\geq1}nG_K(K)\qquad .
\end{equation*}

Tenemos ahora el siguiente resultado.

\medskip
\noindent
\underline{Lema 14.3.4.} \underline{Sea} $V$ \underline{un anillo de
valoración discreta cuyo cuerpo residual} $k$ \underline{tenga
característica} $p>0$ \underline{y sea de tipo finito sobre su cuerpo
primo, y sea} $G_K$ \underline{un esquema en grupos de tipo finito
sobre} $K$, \underline{de rango unipotente cero si}
$\operatorname{car}K=0$. \underline{Entonces el grupo} $g$
\underline{de los elementos infinitamente divisibles de} $G_K(K)$
\underline{es cero; es decir,} (\underline{si} $G_K$ \underline{es
divisible}) \underline{el homomorfismo} (14.3.3) \underline{es
inyectivo}.

Después de efectuar una extensión finita $K'/K$ y reemplazar $V$ por
su normalización en $K'$, podemos suponer que $G_K$ admite una serie de
composición cuyos cocientes sucesivos son de uno de los tipos
siguientes: un grupo finito constante, una variedad abeliana,
$\mathbb G_{mK}$ o $\mathbb G_{aK}$. Nos reducimos, pues,
inmediatamente al caso en que el propio $\sisi{\mathbb G_K}{G_K}$ es
de uno de estos tipos. El caso de un grupo finito constante, o de
$\mathbb G_{aK}$, es trivial. En el caso abeliano, $G_K$ puede
considerarse como la fibra genérica de su modelo de Néron $G$, de modo
que nos reducimos a demostrar que, si $G$ es un esquema de tipo finito
sobre $S$, el subgrupo de $G(S)$ formado por los puntos infinitamente
divisibles es cero.

%%%%%%%%%%  scan idx 527  |  folio 516  %%%%%%%%%%

Esto es cierto para el grupo $G(k)$, como se deduce inmediatamente del
dévissage habitual de $G_o$ y del \underline{teorema de Mordell--Weil}
[16] (que dice que $G_o(k)$ es un $\mathbb Z$-módulo de tipo finito si
$G_o$ es un esquema abeliano sobre el cuerpo de tipo finito $k$). Por
consiguiente, $g$ está contenido en el núcleo $G(S)^{(1)}$ de
$G(S)\longrightarrow G(k)$. El núcleo de
$G(S)^{(1)}\longrightarrow G(V/\underline m^{n+1})$ posee una
filtración cuyos cocientes son isomorfos a grupos vectoriales sobre $k$
(después de suponer, como podemos, que $S$ es completo), y estos son
aniquilados por $p$. Se deduce que $g$ está contenido en el núcleo del
homomorfismo precedente para todo $n$, y por tanto que es cero. El caso
$G_K=\mathbb G_{mK}$ se reduce al precedente, pues
$G_K(K)\sisi{\longrightarrow}{=}K^*$ es una extensión de $\mathbb Z$
(que no tiene elementos infinitamente divisibles no nulos) por
$V^*=\mathbb G_m(S)$, y este último satisface la misma condición por lo
que acabamos de demostrar.

\medskip
\noindent
\underline{Observaciones 14.3.4.1.} Obsérvese que en 14.3.4 es esencial
exigir la divisibilidad por \underline{todos} los enteros $\geq1$, y no
solo por las potencias de un primo fijo $\ell$; equivalentemente, hay
que considerar el núcleo de (14.3.3), en vez de restringir el miembro
de la derecha a un único factor $H^1(K,T_\ell(G_K))$. En efecto, si
$G_K$ proviene de un esquema en grupos liso $G$ sobre $S$ y $S$ es
henseliano, ¡todo elemento del núcleo de
$G(S)\longrightarrow G(k)$ es infinitamente $\ell$-divisible para todo
$\ell\neq p$! El argumento que demuestra 14.3.4 muestra que el subgrupo
de $G_K(K)$ formado por los elementos infinitamente $p$-divisibles (es
decir, el núcleo de
$G_K(K)\longrightarrow H^1(K,T_p(G_K))$ cuando $G_K$ es $p$-divisible)
es un grupo \underline{finito}.

%%%%%%%%%%  scan idx 528  |  folio 517  %%%%%%%%%%

\medskip
\noindent
14.3.5. Las consideraciones precedentes muestran que, cuando el cuerpo
residual $k$ tiene característica $p>0$ y es de tipo finito en sentido
absoluto, el homomorfismo $i$ (14.1.5) está completamente determinado
por su imagen mediante (14.3.1), es decir, por las extensiones
(14.2.2). Para deducir de ello un principio que caracterice $i$ en el
caso general, sin restricción sobre $k$, obsérvese que en el párrafo 7
pudimos construir el esquema $G$ sobre $S$ bajo condiciones
considerablemente más generales que aquellas en las que $S$ es un
rasgo completo y se dispone de una variedad abeliana con reducción
semiestable sobre su cuerpo de fracciones.

Más generalmente, sea $S$ el espectro de un anillo noetheriano normal
$V$, separado y completo para una topología $\underline m$-ádica, y
consideremos un esquema en grupos liso y cuasiproyectivo $A$ sobre $S$,
de fibras conexas. Supongamos que
$A_o=A\times_S S_o$, donde
$S_o=\operatorname{Spec}(V/\underline m)$, es una extensión de un
esquema abeliano $B_o$ por un toro isotrivial $T_o$, y que $A_U$ es un
esquema abeliano, donde $U\subset S-S_o$ es un abierto no vacío. En
$\sisi{7.22}{7.2.2}$ construimos entonces un esquema $G$ sobre $S$, una
extensión de un esquema abeliano $B$ por un toro $T$, caracterizado
salvo isomorfismo único por el isomorfismo (14.1.8)

\begin{equation}\tag{14.3.5.1}
 \widehat G\xrightarrow{\ \sim\ }\widehat A
\end{equation}

\noindent
sobre los completados formales a lo largo de $G_o$ y $A_o$.

%%%%%%%%%%  scan idx 529  |  folio 518  %%%%%%%%%%

Para todo número primo $\ell$ se deduce un isomorfismo (7.3.5)

\begin{equation}\tag{14.3.5.2}
 T_\ell(G)\xrightarrow{\ \sim\ }T_\ell(A)^f\qquad ,
\end{equation}

\noindent
y por tanto, restringiendo a $U$, la estructura de una extensión de
grupos de Barsotti--Tate pro-$\ell$

\begin{equation}\tag{14.3.5.3}
\begin{tikzcd}[row sep=large,column sep=large]
0 \arrow[r] & T_\ell(G_U) \arrow[r] \arrow[d,no head,"\wr"']
 & T_\ell(A_U) \arrow[r]
 & T_\ell(A_U)/T_\ell(A_U)^f \arrow[r] & 0\quad .\\
& T_\ell(A_U)^f
\end{tikzcd}
\end{equation}

Supongamos ahora dado un segundo esquema en grupos $A'$ sobre $S$ que
satisface las mismas hipótesis que $A$, junto con una biextensión $W$
de $(A,A')$ por $\mathbb G_m$ que sobre $U$ induce una dualidad entre
los esquemas abelianos $A_U$ y $A'_U$. (De hecho, por VIII 7.1, dar $W$
equivale a dar tal dualidad entre $A_U$ y $A'_U$.) Entonces $W$ define
emparejamientos

\begin{equation}\tag{14.3.5.4}
 T_\ell(A_U)\times T_\ell(A'_U)\longrightarrow\mathbb Z_\ell(1)
 \qquad ,
\end{equation}

\noindent
que dan lugar al teorema de ortogonalidad (7.4.8):
$T_\ell(A_U)^f$ es el ortogonal de
$T_\ell(A'_U)^t=T_\ell(T'_U)$, donde $T'$ es la parte tórica de la
extensión de Raynaud $G'={A'}^{\natural}$ asociada a $A'$. Se obtiene
así de nuevo un isomorfismo

\begin{equation}\tag{14.3.5.5}
 T_\ell(A_U)/T_\ell(A_U)^f
 \simeq\underline M_U\otimes\mathbb Z_\ell\qquad ,
\end{equation}

\noindent
donde $\underline M=D(T')$ es el retículo sobre $S$ dual del toro
$T'$. Por consiguiente, $\sisi{(14.3.8)}{(14.3.5.3)}$ y
$\sisi{(14.3.10)}{(14.3.5.5)}$ dan una extensión

\begin{equation}\tag{14.3.5.6}
 0\longrightarrow T_\ell(G_U)\longrightarrow T_\ell(A_U)
 \longrightarrow\underline M_U\otimes\mathbb Z_\ell
 \longrightarrow0\qquad .
\end{equation}

Entendido esto, podemos enunciar la siguiente conjetura.
